\begin{lemma}[Lemma 2.5]
Let $n\ge2$, $\gamma>0$ and $\lambda_1\in(0,1)$, and suppose the Sobolev--Poincar\'e inequality on balls holds in dimension $n$
(\noderef{CitedSobolevPoincareBall}).  Then there exists
$\delta=\delta(n,\gamma,\lambda_1)\in(0,1)$ with the following property.  Let
$x_0\in\R^n$, $R>0$, let $f\in W^{1,2}(B_R(x_0))$ with weak gradient $g$
(\noderef{MemW1pOn}), suppose $f$ is essentially bounded on $B_R(x_0)$ and put
$M=\esssup_{B_R(x_0)}f$, $m=\essinf_{B_R(x_0)}f$.  Suppose that
\begin{enumerate}
\item the De Giorgi level bound (2.1) of the thesis with constant $\gamma$
  (\noderef{DeGiorgiLevelBound}) holds for $f$ at every level
  $\ell\in\bigl[\frac{m+M}2,M\bigr]$ and all radii
  $\rho_0,\rho_1\in[\lambda_1R,R]$ with $\rho_0<\rho_1$; and
\item condition (2.11) of the thesis,
\[
\int_{B_{\rho_0}}|\nabla(\ell-f)^+|^2dx\le\frac{\gamma}{(\rho_1-\rho_0)^2}
\max_{B_{\rho_1}}[(\ell-f)^+]^2\bigl|\{x\in B_{\rho_1}:f(x)<\ell\}\bigr| ,
\]
holds for every $\ell\in\bigl[m,\frac{m+M}2\bigr]$ and all
$\rho_0,\rho_1\in[\lambda_1R,R]$ with $\rho_0<\rho_1$.
\end{enumerate}
Then
\[
\operatorname*{osc}(f,B_{\lambda_1R}(x_0))\le\delta\,\operatorname*{osc}(f,B_R(x_0)) .
\]

The hypothesis \noderef{CitedSobolevPoincareBall} is the cited result (C6) of the
paper's Note for the formalization, carried here as an explicit argument because
the proof uses it --- it enters through the constants $\theta_0$ of \noderef{LevelTruncationMaximumPrinciple} and $\lambda_2$ of \noderef{LevelSetMeasureImprovement} --- which is what the
Note permits for the results the thesis quotes without proof.  Although it is
bound before the existential for $\delta$, it introduces no dependence of
$\delta$ on any quantity outside the list the paper names: (C6) is a statement
about $n$ alone, and $n$ is already bound first, so the hypothesis carries no
numerical content of its own.
\end{lemma}
\begin{proof}
\emph{Step 0: the choice of $\delta$.}  Since $\lambda_1\in(0,1)$ we have
$\sqrt{\lambda_1}\in(\lambda_1,1)$, so $\sqrt{\lambda_1}\in(0,1)$.  Let
$\theta_0=\theta_0\bigl(n,\gamma,\tfrac12,\sqrt{\lambda_1}\bigr)>0$ be the
constant furnished by \noderef{LevelTruncationMaximumPrinciple} --- whose
statement requires the hypothesis \noderef{CitedSobolevPoincareBall}, supplied
here by the identical hypothesis of the present lemma, and likewise for
\noderef{LevelSetMeasureImprovement} below --- for the
parameters $n$, $\gamma$, $\lambda_0=\tfrac12$ and $\sqrt{\lambda_1}$, put
\[
\theta_1=\min\bigl\{\theta_0,\tfrac12\bigr\}\in(0,1),
\]
and let $\lambda_2=\lambda_2\bigl(n,\gamma,\sqrt{\lambda_1},\tfrac12,\theta_1\bigr)
\in(0,1)$ be the constant furnished by \noderef{LevelSetMeasureImprovement} for
the parameters $n$, $\gamma$, $\sqrt{\lambda_1}$, $c_0=\tfrac12$ and $\theta_1$.
Define
\[
\delta=1-\frac{\lambda_2}4 .
\]
Since $\lambda_2\in(0,1)$ we have $\frac{\lambda_2}4\in(0,\tfrac14)$, hence
$\delta\in(\tfrac34,1)\subseteq(0,1)$ as the conclusion requires.  All four of
$\theta_0,\theta_1,\lambda_2,\delta$ were constructed from $n,\gamma,\lambda_1$
alone, before any of $x_0,R,m,M,f,g$, so $\delta=\delta(n,\gamma,\lambda_1)$.

Now fix $x_0,R,m,M,f,g$ satisfying the hypotheses.  Write
$A^+_{\sigma,\rho}(v)=\{x\in B_\rho(x_0):v(x)>\sigma\}$ for a function $v$, and
abbreviate $\operatorname*{osc}(v,S)=\esssup_Sv-\essinf_Sv$.

\emph{Step 1: the hypothesis package, and two symmetries.}  We record three
facts.  The first, (1b), is used at three later places and is proved once here:
hypothesis (2) of the lemma is a bare \emph{display}, while hypothesis (1) is the
two-clause predicate \noderef{DeGiorgiLevelBound}, and the identification of the
one with the other at the reflected pair is the paper's ``it is clear now'' at
l.~496--497.  It is clear in the paper because there (2.1) is also a bare display
and ``max means essential max'' is a convention announced at l.~290--291; here
that convention has been promoted into clause (1) of the predicate --- it has to
be, since $\esssup$ of a function that is not essentially bounded above is a junk
value --- and the promotion is exactly what turns the paper's one-line remark
into a step owing a derivation.

\emph{(1a) Translation.}  Let $c\in\R$ and put $\tilde f=f-c$.  Then
$\esssup_S\tilde f=\esssup_Sf-c$ and $\essinf_S\tilde f=\essinf_Sf-c$ for every
$S$ of positive measure, so $\operatorname*{osc}(\tilde f,S)=\operatorname*{osc}(f,S)$;
in particular the conclusion for $\tilde f$ and the conclusion for $f$ are the
same inequality.  Also $\tilde f\in W^{1,2}(B_R)$ with weak gradient $g$: this is
clause (2) of \noderef{WeakGradientLinear} on the open set $B_R$ with $q=2$,
$a=b=1$, $f_1=f$, $g_1=g$, and $f_2$ the constant function $x\mapsto-c$ with
$g_2=0$, whose own $W^{1,2}(B_R)$ membership is clause (4) of the same node
(applicable because $|B_R|<\infty$); the combination $af_1+bf_2$ is $f-c=\tilde f$
and its weak gradient $ag_1+bg_2$ is $g+0=g$.  And $\tilde f$
is essentially bounded on $B_R$ iff $f$ is.  Finally, for every
$\ell,\rho_0,\rho_1$, the De Giorgi level bound
(\noderef{DeGiorgiLevelBound}) with constant $\gamma$ holds for $(f,g)$ at level
$\ell$ and radii $\rho_0,\rho_1$ if and only if it holds for $(\tilde f,g)$ at
level $\ell-c$ and the same radii: indeed
$\tilde f-(\ell-c)=f-\ell$, so the truncations $(\tilde f-(\ell-c))^+$ and
$(f-\ell)^+$ coincide, the level sets
$\{\tilde f>\ell-c\}$ and $\{f>\ell\}$ coincide, and the weak gradient is
unchanged; every ingredient of that predicate --- the essential boundedness of
the truncation, the integral of $|g|^2$ over
$B_{\rho_0}\cap\{\,\cdot>\text{level}\}$, the essential supremum of the squared
truncation over $B_{\rho_1}$, and the measure of
$B_{\rho_1}\cap\{\,\cdot>\text{level}\}$ --- is therefore literally the same
number for the two.

\emph{(1b) The display (2.11) of the thesis at a level $\ell$ is the level bound for
$(-f,-g)$ at the level $-\ell$.}  Fix $\ell\in\R$ and radii
$\rho_0,\rho_1\in[\lambda_1R,R]$ with $\rho_0<\rho_1$; in particular
$\rho_1\le R$, so $B_{\rho_1}\subseteq B_R$.  By (1a)'s computation --- or
directly by clause (2) of \noderef{WeakGradientLinear} on $B_R$ with $q=2$,
$a=-1$, $b=0$ and both pairs $(f,g)$ --- the pair $(-f,-g)$ satisfies
$-f\in W^{1,2}(B_R)$ with weak gradient $-g$ (\noderef{MemW1pOn}).  We claim
that the two clauses of
$\mathrm{DeGiorgiLevelBound}(n,\gamma,-f,-g,x_0,-\ell,\rho_0,\rho_1)$
(\noderef{DeGiorgiLevelBound}) hold if and only if the display of hypothesis (2)
holds at $\ell,\rho_0,\rho_1$, and that the two are in fact the \emph{same}
inequality once clause (1) is in hand.

\emph{Clause (1) is free, from the lemma's own hypothesis.}  The lemma assumes
$f$ essentially bounded on $B_R$; fix $C\ge0$ with $|f|\le C$ almost everywhere
on $B_R$.  The clause asks for an essential bound on
$\bigl((-f)-(-\ell)\bigr)^+=(\ell-f)^+$ over $B_{\rho_1}$, and
$B_{\rho_1}\subseteq B_R$, so almost everywhere on $B_{\rho_1}$ either
$\ell-f\le0$, in which case $(\ell-f)^+=0\le|\ell|+C$, or $\ell-f>0$, in which
case $(\ell-f)^+=\ell-f\le|\ell|+|f|\le|\ell|+C$.  So $|\ell|+C$ serves, and
clause (1) holds.  This is precisely the argument the definition node points at:
\noderef{DeGiorgiLevelBound} says of its clause (1) that ``it is part of the
hypothesis package, exactly as in the paper, where $(f-\ell)^+$ is bounded on
$B_{\rho_1}$ because $f$ is''.  The same computation with $f$ and $\ell$ in place
of $-f$ and $-\ell$ gives clause (1) for the unreflected pair, which is what the
other two sites below need.

\emph{Clause (2) is the display, term by term.}  Three identifications, none of
them an inequality.
\begin{itemize}
\item \emph{The integrand.}  Apply \noderef{WeakGradientPositivePart} with
  $q=2$, the level $c=-\ell$, the open set $B_R$ --- whose volume is finite
  (Preamble: the volume of a ball) --- and the pair $(-f,-g)$: it gives
  $\bigl((-f)-(-\ell)\bigr)^+=(\ell-f)^+\in W^{1,2}(B_R)$ with weak gradient
  $x\mapsto-g(x)$ on $\{-\ell<-f\}$ and $0$ elsewhere.  A weak gradient is
  determined almost everywhere (Preamble: the fundamental lemma of the calculus
  of variations, second consequence), so the $\nabla(\ell-f)^+$ of the display
  \emph{is} that field almost everywhere, and pointwise almost everywhere
  \[
  \bigl|\nabla(\ell-f)^+\bigr|^2=\|{-g}\|^2\chi_{\{-\ell<-f\}}
  =\|g\|^2\chi_{\{f<\ell\}} ,
  \]
  using $\|-g\|=\|g\|$ and $\{-\ell<-f\}=\{f<\ell\}$.  Hence the display's
  integral over all of $B_{\rho_0}$ equals the predicate's integral of
  $\|-g\|^2$ over $B_{\rho_0}\cap\{-f>-\ell\}$: the integrand vanishes off that
  set, which is $B_{\rho_0}\cap\{f<\ell\}$.
\item \emph{The essential-supremum factor.}  $\bigl((-f)-(-\ell)\bigr)^+
  =(\ell-f)^+$ pointwise, so
  $\bigl(\esssup_{B_{\rho_1}}((-f)-(-\ell))^+\bigr)^2$ and
  $\max_{B_{\rho_1}}[(\ell-f)^+]^2$ are the same number --- the essential
  supremum being the paper's ``max'' by the convention at l.~290--291, and being
  finite by clause (1) just proved, which is what stops it being a junk value.
\item \emph{The measure factor.}  $\{x\in B_{\rho_1}:-f(x)>-\ell\}
  =\{x\in B_{\rho_1}:f(x)<\ell\}$, so the two measures are the measure of one
  and the same set.
\end{itemize}
Both sides of clause (2) are therefore literally the two sides of the display, so
each implies the other.  Together with clause (1) this proves the claim.  Read in
the other direction --- with $f$, $\ell$ in place of $-f$, $-\ell$ --- the same
three identifications show that the predicate for $(f,g)$ at level $\ell$ is the
display (2.1) of the thesis at $\ell$, which is the form hypothesis (1) is quoted in.

\emph{(1c) Reflection.}  Put $\hat f=-f$.  Then $\hat f\in W^{1,2}(B_R)$ with weak
gradient $-g$, again by clause (2) of \noderef{WeakGradientLinear} on $B_R$ with
$q=2$, this time with $a=-1$, $b=0$ and both pairs taken to be
$(f_1,g_1)=(f_2,g_2)=(f,g)$: the combination is $-f+0\cdot f=\hat f$ and its weak
gradient is $(-1)\cdot g+0\cdot g=-g$.  Then
$\esssup_S\hat f=-\essinf_Sf$ and $\essinf_S\hat f=-\esssup_Sf$, so
$\operatorname*{osc}(\hat f,S)=\operatorname*{osc}(f,S)$, and $\hat f$ is
essentially bounded on $B_R$ with $\esssup_{B_R}\hat f=-m$,
$\essinf_{B_R}\hat f=-M$.  Moreover the hypothesis package is invariant under
$f\mapsto\hat f$: by (1b), hypothesis (2) of the lemma at a level $\ell$ is
exactly the level bound with constant $\gamma$ for
$(-f,-g)=(\hat f,-g)$ at level $-\ell$, so hypothesis (2) gives that bound for
every $\ell\in[m,\frac{m+M}2]$, i.e.\ at every level in
$[-\frac{m+M}2,-m]=[\frac{\hat m+\hat M}2,\hat M]$ where $\hat M=-m$ and
$\hat m=-M$; so hypothesis (2) for $f$ is hypothesis (1) for $\hat f$.
Symmetrically, hypothesis (1) for $f$ is the level bound for $(f,g)=
(-\hat f,-(-g))$ at every level in $[\frac{m+M}2,M]=[-\frac{\hat m+\hat M}2,-\hat m]$,
and (1b) read at the pair $(\hat f,-g)$ and the level $-\ell$ --- for which
$(-\hat f,-(-g))=(f,g)$ and $-(-\ell)=\ell$ --- turns that level bound into the
display (2.11) of the thesis for $\hat f$ at the level $-\ell$; so it is hypothesis
(2) for $\hat f$.  Clause (1) of each of these predicates is supplied by the
computation in (1b), applied with the essential bound $C$ for $f$, which is also
one for $\hat f$.  Hence $\hat f$ satisfies both
hypotheses with the same $\gamma$, $\lambda_1$, $x_0$ and $R$.

\emph{Step 2: reduction to $m=-M$, and the case $M=0$.}  By (1a) applied with
$c=\frac{m+M}2$ we may replace $f$ by $f-\frac{m+M}2$: the conclusion is
unchanged, and the hypotheses are preserved.  For hypothesis (1) this is (1a)
directly, since that hypothesis is the predicate.  For hypothesis (2) it is
(1b) first and then (1a): by (1b), hypothesis (2) at a level $\ell$ is the level
bound for the pair $(-f,-g)$ at level $-\ell$, and (1a) applied to that pair with
$c'=-\frac{m+M}2$ --- replacing $-f$ by $-f+\frac{m+M}2=-(f-\frac{m+M}2)$ and
shifting its level from $-\ell$ to $-\ell-c'=-\bigl(\ell-\frac{m+M}2\bigr)$ ---
gives the level bound for the reflected translate at the translated level, which
by (1b) read backwards is hypothesis (2) for $f-\frac{m+M}2$ at the level
$\ell-\frac{m+M}2$.  So both hypotheses survive, with the level ranges
$[\frac{m+M}2,M]$ and $[m,\frac{m+M}2]$ translated to
$[0,\frac{M-m}2]$ and $[-\frac{M-m}2,0]$ respectively, which are exactly the
ranges $[\frac{m'+M'}2,M']$ and $[m',\frac{m'+M'}2]$ for the new
$M'=\frac{M-m}2$ and $m'=-\frac{M-m}2$.  So we may and do assume
\[
m=-M,\qquad\text{hence}\qquad M\ge0
\]
(the latter because $m\le M$, i.e.\ $-M\le M$), and then
$\frac{m+M}2=0$, so hypothesis (1) holds at every level in $[0,M]$ and
hypothesis (2) at every level in $[-M,0]$, and
$\operatorname*{osc}(f,B_R)=M-(-M)=2M$.

If $M=0$ then $m=M=0$, so $\operatorname*{osc}(f,B_R)=0$ and, since
$B_{\lambda_1R}\subseteq B_R$ gives
$\essinf_{B_R}f\le\essinf_{B_{\lambda_1R}}f\le\esssup_{B_{\lambda_1R}}f
\le\esssup_{B_R}f$, also $\operatorname*{osc}(f,B_{\lambda_1R})=0$; the conclusion
$0\le\delta\cdot0$ holds.  Assume from now on $M>0$.

\emph{Step 3: reduction to $|A^+_{0,\sqrt{\lambda_1}R}(f)|\le\frac12
|B_{\sqrt{\lambda_1}R}|$.}  The sets $A^+_{0,\sqrt{\lambda_1}R}(f)=
\{x\in B_{\sqrt{\lambda_1}R}:f(x)>0\}$ and
$A^+_{0,\sqrt{\lambda_1}R}(-f)=\{x\in B_{\sqrt{\lambda_1}R}:f(x)<0\}$ are
disjoint subsets of $B_{\sqrt{\lambda_1}R}$, so their measures sum to at most
$|B_{\sqrt{\lambda_1}R}|$ and therefore at least one of them is at most
$\frac12|B_{\sqrt{\lambda_1}R}|$.  If it is the first, proceed.  If it is the
second, replace $f$ by $\hat f=-f$: by (1c) --- whose invariance of the
hypothesis package rests on (1b) at both levels --- this preserves both
hypotheses, the
value of $M$ (for $\hat f$ one has $\esssup_{B_R}\hat f=-m=M$ and
$\essinf_{B_R}\hat f=-M$, so $\hat f$ again satisfies $\hat m=-\hat M$ with
$\hat M=M$), and the conclusion, while making the required inequality hold.  So
we may and do assume
\[
\bigl|A^+_{0,\sqrt{\lambda_1}R}(f)\bigr|\le\tfrac12\bigl|B_{\sqrt{\lambda_1}R}\bigr| .
\]

\emph{Step 4: improving the measure of a level set, via Lemma 2.4.}  We apply
\noderef{LevelSetMeasureImprovement} with parameters $n$, $\gamma$,
$\sqrt{\lambda_1}$ in place of its $\lambda_1$, $c_0=\frac12$ and $\theta_1$, on
the ball $B_R(x_0)$, at the level $k=0$.  Its hypotheses hold:
$f\in W^{1,2}(B_R)$ with weak gradient $g$ and $f$ essentially bounded above on
$B_R$ with $M=\esssup_{B_R}f$ are hypotheses of the present lemma;
$k=0<M$ by Step 2; the level bound with constant $\gamma$ at
$\rho_0=\sqrt{\lambda_1}R$, $\rho_1=R$ and every level $\ell\in[0,M]$ is
supplied by hypothesis (1), because $\sqrt{\lambda_1}R\in[\lambda_1R,R]$ and
$\sqrt{\lambda_1}R<R$ (as $\lambda_1<\sqrt{\lambda_1}<1$); and
$\bigl|A^+_{0,\sqrt{\lambda_1}R}\bigr|\le\frac12\bigl|B_{\sqrt{\lambda_1}R}\bigr|$
is Step 3.  Since the $\lambda_2$ of Step 0 is the constant that
\noderef{LevelSetMeasureImprovement} produces for exactly these parameters, its
conclusion with $k=0$, namely
$\lambda_2\cdot0+(1-\lambda_2)M=(1-\lambda_2)M$, reads
\[
\bigl|\{x\in B_{\sqrt{\lambda_1}R}(x_0):f(x)>(1-\lambda_2)M\}\bigr|
\le\theta_1\bigl|B_{\sqrt{\lambda_1}R}\bigr| .
\tag{A}
\]

\emph{Step 5: the two cases.}  Since $f$ is essentially bounded on $B_R$ and
$B_{\sqrt{\lambda_1}R}\subseteq B_R$ has positive measure, $f$ is essentially
bounded on $B_{\sqrt{\lambda_1}R}$; put
\[
M'=\esssup_{B_{\sqrt{\lambda_1}R}}f ,\qquad\text{so}\qquad M'\le M
\]
by monotonicity of the essential supremum under
$B_{\sqrt{\lambda_1}R}\subseteq B_R$.  Exactly one of the following holds.

\emph{Case 1: $(1-\lambda_2)M<M'$.}  We apply
\noderef{LevelTruncationMaximumPrinciple} with parameters $n$, $\gamma$,
$\lambda_0=\frac12$ and $\sqrt{\lambda_1}$ in place of its $\lambda_1$, on the
ball $B_{\sqrt{\lambda_1}R}(x_0)$ --- that is, with its $R$ equal to
$\sqrt{\lambda_1}R$ --- at the level $k=(1-\lambda_2)M$.  Its hypotheses hold:
$\sqrt{\lambda_1}R>0$; $f\in W^{1,2}(B_{\sqrt{\lambda_1}R})$ with weak gradient
$g$ --- the hypothesis of the present lemma gives
$f\in W^{1,2}(B_R)$ with weak gradient $g$ (\noderef{MemW1pOn}), and clause (1) of
\noderef{WeakGradientRestrict}, applied with $q=2$, $\Omega=B_R(x_0)$ and
$\Omega'=B_{\sqrt{\lambda_1}R}(x_0)$ --- both open, and
$B_{\sqrt{\lambda_1}R}\subseteq B_R$ because $\sqrt{\lambda_1}<1$ --- transports it
to the smaller ball with the same weak gradient; $f$ is essentially bounded above on
$B_{\sqrt{\lambda_1}R}$ with
essential supremum $M'$, just established; $k=(1-\lambda_2)M<M'$, which is the
case assumption; the level bound with constant $\gamma$ holds at every level
$\ell\in[k,\frac12k+\frac12M']$ and all radii
$\rho_0,\rho_1\in[\sqrt{\lambda_1}\cdot\sqrt{\lambda_1}R,\sqrt{\lambda_1}R]
=[\lambda_1R,\sqrt{\lambda_1}R]$ with $\rho_0<\rho_1$, because those radii lie in
$[\lambda_1R,R]$ and those levels lie in $[0,M]$ --- the lower end because
$k=(1-\lambda_2)M\ge0$ (as $\lambda_2\le1$ and $M\ge0$) and the upper because
$\frac12k+\frac12M'\le\frac12M+\frac12M=M$ (as $k\le M$ and $M'\le M$) --- so
hypothesis (1) supplies them; and the measure hypothesis
$|\{x\in B_{\sqrt{\lambda_1}R}:f>k\}|\le\theta_0|B_{\sqrt{\lambda_1}R}|$ follows
from (A) together with $\theta_1\le\theta_0$.  Its conclusion, with
$\lambda_0=\frac12$ and its $\lambda_1$ equal to $\sqrt{\lambda_1}$, is that
\[
f\le\tfrac12(1-\lambda_2)M+\tfrac12M'
\qquad\text{a.e.\ in }B_{\sqrt{\lambda_1}\cdot\sqrt{\lambda_1}R}(x_0)
=B_{\lambda_1R}(x_0) .
\]
Since $M'\le M$, the right-hand side is at most
$\frac12(1-\lambda_2)M+\frac12M=M-\frac{\lambda_2}2M$, so
\[
f\le\Bigl(1-\frac{\lambda_2}2\Bigr)M\qquad\text{a.e.\ in }B_{\lambda_1R}(x_0) .
\tag{B}
\]
Note that only $M'\le M$ was needed here.

\emph{Case 2: $(1-\lambda_2)M\ge M'$.}  Here
\noderef{LevelTruncationMaximumPrinciple} is not needed at all.  By the
definition of $M'$ we have $f\le M'$ a.e.\ in $B_{\sqrt{\lambda_1}R}$, hence a
fortiori a.e.\ in $B_{\lambda_1R}\subseteq B_{\sqrt{\lambda_1}R}$ (the inclusion
because $\lambda_1<\sqrt{\lambda_1}$).  Therefore, a.e.\ in $B_{\lambda_1R}$,
\[
f\le M'\le(1-\lambda_2)M\le\Bigl(1-\frac{\lambda_2}2\Bigr)M ,
\]
the last step because $(1-\frac{\lambda_2}2)M-(1-\lambda_2)M
=\frac{\lambda_2}2M\ge0$, using only $\lambda_2\ge0$ and $M\ge0$.  So (B) holds
in this case too.

\emph{Step 6: from the one-sided bound to the oscillation.}  By (B),
\[
\esssup_{B_{\lambda_1R}}f\le\Bigl(1-\frac{\lambda_2}2\Bigr)M .
\]
For the lower bound we use nothing new: $B_{\lambda_1R}\subseteq B_R$ gives
\[
\essinf_{B_{\lambda_1R}}f\ge\essinf_{B_R}f=m=-M .
\]
Hence
\[
\operatorname*{osc}(f,B_{\lambda_1R})
=\esssup_{B_{\lambda_1R}}f-\essinf_{B_{\lambda_1R}}f
\le\Bigl(1-\frac{\lambda_2}2\Bigr)M-(-M)
=\Bigl(2-\frac{\lambda_2}2\Bigr)M
=\Bigl(1-\frac{\lambda_2}4\Bigr)\cdot2M ,
\]
the last equality being the identity
$2-\frac{\lambda_2}2=2\bigl(1-\frac{\lambda_2}4\bigr)$.  Since
$\operatorname*{osc}(f,B_R)=2M$ by Step 2, this is
\[
\operatorname*{osc}(f,B_{\lambda_1R}(x_0))
\le\Bigl(1-\frac{\lambda_2}4\Bigr)\operatorname*{osc}(f,B_R(x_0))
=\delta\,\operatorname*{osc}(f,B_R(x_0)) ,
\]
which is the assertion.  The halving of the rate is worth noting: the gain
$\frac{\lambda_2}2M$ was obtained on the upper bound only, while the lower bound
$-M$ was carried over unchanged, so the improvement is measured against the full
oscillation $2M$ rather than against $M$.  That is exactly why
$\delta=1-\frac{\lambda_2}4$ and not $1-\frac{\lambda_2}2$.
\end{proof}
